\documentclass[aspectratio=169,10pt]{beamer}
\usetheme[numbering=fraction,progressbar=frametitle]{metropolis}
\usepackage{fontspec}
\setsansfont{Latin Modern Sans}
\usepackage{amsmath,amssymb,mathtools,graphicx,booktabs}
\usepackage{polyglossia}
\setdefaultlanguage{french}
\setbeamertemplate{navigation symbols}{}
\setbeamerfont{frametitle}{size=\large}
\setbeamerfont{normal text}{size=\normalsize}
\setlength{\parskip}{.45em}
\hypersetup{pdfauthor={},pdfsubject={Lecture slides},pdfkeywords={},pdftitle={Fonctions exponentielles}}
\providecommand{\N}{\mathbb{N}}
\providecommand{\Z}{\mathbb{Z}}
\providecommand{\D}{\mathbb{D}}
\providecommand{\Q}{\mathbb{Q}}
\providecommand{\R}{\mathbb{R}}
\providecommand{\C}{\mathbb{C}}
\providecommand{\vect}[1]{\overrightarrow{#1}}
\providecommand{\vi}{\vec{\imath}}
\providecommand{\vj}{\vec{\jmath}}
\providecommand{\vk}{\vec{k}}
\providecommand{\norm}[1]{\left\lVert #1\right\rVert}
\providecommand{\vc}[2]{\begin{pmatrix}#1\\#2\end{pmatrix}}
\providecommand{\deter}[2]{\det\left(#1,#2\right)}
\providecommand{\abs}[1]{\left\lvert #1\right\rvert}
\providecommand{\dd}{\,\mathrm{d}}
\providecommand{\e}{\mathrm{e}}
\providecommand{\Card}{\operatorname{Card}}
\providecommand{\E}{\operatorname{E}}
\providecommand{\V}{\operatorname{V}}
\providecommand{\degre}{^{\circ}}
\providecommand{\un}[1]{\,\mathrm{#1}}
\providecommand{\sci}[2]{#1\times10^{#2}}
\providecommand{\pgcd}{\operatorname{PGCD}}
\providecommand{\ppcm}{\operatorname{PPCM}}
\title{Fonctions exponentielles}
\subtitle{Cours et exemples}
\date{}
\author{}
\begin{document}
\begin{frame}[plain]
\vfill
{\usebeamerfont{title}\inserttitle\par}
\medskip
{\usebeamerfont{subtitle}\insertsubtitle\par}
\medskip
{\usebeamercolor[fg]{progress bar}\rule{\textwidth}{.5pt}}
\vfill
\end{frame}
\begin{frame}{Cours de référence}
Si \(a>1\), \(x\mapsto a^x\) est strictement croissante. Si \(0<a<1\), elle est strictement décroissante. Pour \(f(x)=ka^x\), multiplier par \(k<0\) inverse le sens de variation. Dans tous les cas, \(a^x>0\).
\end{frame}
\begin{frame}{Cours de référence}
En général, \(a^x+a^y\neq a^{x+y}\). Les propriétés portent sur les produits et les quotients, pas sur les sommes.
\end{frame}
\begin{frame}{Cours de référence}
Un chiffre d'affaires passe de \(80\,000\) à \(95\,000\) euros en quatre ans. Le coefficient global vaut \(Q=\dfrac{95\,000}{80\,000}=1{,}1875\). Le coefficient annuel moyen est \(q=1{,}1875^{\dfrac14}\approx1{,}0439\) : le taux annuel moyen est donc environ \(4{,}39\,\%\).
\end{frame}
\begin{frame}{Définition et sens de variation de \(a^x\)}
\begin{itemize}
\item Donner du sens à \(a^x\).
\item Relier base et variation.
\item Lire des courbes.

\end{itemize}
\end{frame}
\begin{frame}{Définition et sens de variation de \(a^x\)}
\begin{columns}[T,onlytextwidth]
\begin{column}{.57\textwidth}
À l'aide d'une calculatrice, compléter un tableau de valeurs de \(1{,}5^x\) et \(0{,}7^x\) pour \(x=-2,-1,0,1,2,3\). Tracer les points, conjecturer les variations et identifier le point commun aux deux courbes.
\end{column}
\begin{column}{.4\textwidth}
\centering
\includegraphics[width=\linewidth,height=4.5cm,keepaspectratio]{figures/stmg-02-2.pdf}
\end{column}
\end{columns}
\end{frame}
\begin{frame}{Définition et sens de variation de \(a^x\)}
Pour \(a>0\), \(a^x\) est défini pour tout réel \(x\). La fonction est croissante si \(a>1\) et décroissante si \(0<a<1\).
\end{frame}
\begin{frame}{Définition et sens de variation de \(a^x\)}
Reprendre les données de l'activité et rédiger la solution dans l'ordre : définition des objets, propriété utilisée, calcul exact ou approché, puis phrase de réponse. La réponse modèle figure dans l'édition enseignant.
\end{frame}
\begin{frame}{Définition et sens de variation de \(a^x\)}
\textbf{Conduite en 1 h :} 7 min d'automatismes ; 10 min de mise en commun et trace ; 35 min d'exercices essentiels ; 5 min de correction ; 3 min de billet de sortie.

\textbf{Conduite en 2 h :} ajouter 20 min de recherche de l'activité, 20 min d'institutionnalisation, un second exercice autonome et 10 min de correction comparée.

\textbf{Différenciation :} fournir aux élèves fragiles une fiche à quatre cases (données, modèle, calcul, conclusion). Demander aux élèves rapides de produire un contre-exemple, une justification ou une seconde méthode.
\end{frame}
\begin{frame}{Puissances et propriétés algébriques}
\begin{itemize}
\item Transformer des puissances.
\item Utiliser les exposants négatifs.
\item Repérer une règle inapplicable.

\end{itemize}
\end{frame}
\begin{frame}{Puissances et propriétés algébriques}
Simplifier \(2^3\times2^{-5}\), \(5^{x+2}\div5^x\) et \((3^x)^2\). Inventer ensuite un contre-exemple numérique à \(a^x+a^y=a^{x+y}\).
\end{frame}
\begin{frame}{Puissances et propriétés algébriques}
\(a^{x+y}=a^xa^y\), \(a^{x-y}=\dfrac{a^x}{a^y}\) et \((a^x)^y=a^{xy}\). Ces règles exigent une même base positive.
\end{frame}
\begin{frame}{Puissances et propriétés algébriques}
Reprendre les données de l'activité et rédiger la solution dans l'ordre : définition des objets, propriété utilisée, calcul exact ou approché, puis phrase de réponse. La réponse modèle figure dans l'édition enseignant.
\end{frame}
\begin{frame}{Puissances et propriétés algébriques}
\textbf{Conduite en 1 h :} 7 min d'automatismes ; 10 min de mise en commun et trace ; 35 min d'exercices essentiels ; 5 min de correction ; 3 min de billet de sortie.

\textbf{Conduite en 2 h :} ajouter 20 min de recherche de l'activité, 20 min d'institutionnalisation, un second exercice autonome et 10 min de correction comparée.

\textbf{Différenciation :} fournir aux élèves fragiles une fiche à quatre cases (données, modèle, calcul, conclusion). Demander aux élèves rapides de produire un contre-exemple, une justification ou une seconde méthode.
\end{frame}
\begin{frame}{Fonctions \(x\mapsto k a^x\) : variations et graphique}
\begin{itemize}
\item Étudier \(ka^x\).
\item Combiner signe et variation.
\item Relier formule et graphique.

\end{itemize}
\end{frame}
\begin{frame}{Fonctions \(x\mapsto k a^x\) : variations et graphique}
\begin{columns}[T,onlytextwidth]
\begin{column}{.57\textwidth}
Les quatre courbes \(\mathcal C_1\), \(\mathcal C_2\), \(\mathcal C_3\) et \(\mathcal C_4\) ci-dessous représentent, dans un ordre inconnu, les fonctions \(2^x\), \(0{,}6^x\), \(-2^x\) et \(-0{,}6^x\). Associer chaque fonction à sa courbe. Justifier chaque choix par le signe, le point d'ordonnée à l'origine et le sens de variation.
\end{column}
\begin{column}{.4\textwidth}
\centering
\includegraphics[width=\linewidth,height=4.5cm,keepaspectratio]{figures/stmg-02-3.pdf}
\end{column}
\end{columns}
\end{frame}
\begin{frame}{Fonctions \(x\mapsto k a^x\) : variations et graphique}
Pour \(f(x)=ka^x\), le signe est celui de \(k\). Si \(k<0\), la multiplication par \(k\) inverse le sens de variation de \(a^x\).
\end{frame}
\begin{frame}{Fonctions \(x\mapsto k a^x\) : variations et graphique}
Reprendre les données de l'activité et rédiger la solution dans l'ordre : définition des objets, propriété utilisée, calcul exact ou approché, puis phrase de réponse. La réponse modèle figure dans l'édition enseignant.
\end{frame}
\begin{frame}{Fonctions \(x\mapsto k a^x\) : variations et graphique}
\textbf{Conduite en 1 h :} 7 min d'automatismes ; 10 min de mise en commun et trace ; 35 min d'exercices essentiels ; 5 min de correction ; 3 min de billet de sortie.

\textbf{Conduite en 2 h :} ajouter 20 min de recherche de l'activité, 20 min d'institutionnalisation, un second exercice autonome et 10 min de correction comparée.

\textbf{Différenciation :} fournir aux élèves fragiles une fiche à quatre cases (données, modèle, calcul, conclusion). Demander aux élèves rapides de produire un contre-exemple, une justification ou une seconde méthode.
\end{frame}
\begin{frame}{Taux d'évolution moyen}
\begin{itemize}
\item Calculer un taux moyen.
\item Distinguer taux global et moyen.
\item Interpréter un exposant \(\dfrac1n\).

\end{itemize}
\end{frame}
\begin{frame}{Taux d'évolution moyen}
Le prix d'un service passe de \(18\) à \(23{,}50\) euros en cinq ans. Calculer le coefficient global, puis le coefficient annuel moyen et le taux annuel moyen. Vérifier en recomposant cinq fois l'évolution.
\end{frame}
\begin{frame}{Taux d'évolution moyen}
Si \(Q\) est le coefficient global de \(n\) évolutions identiques, le coefficient moyen est \(q=Q^{\dfrac1n}\) et le taux moyen est \(q-1\).
\end{frame}
\begin{frame}{Taux d'évolution moyen}
Reprendre les données de l'activité et rédiger la solution dans l'ordre : définition des objets, propriété utilisée, calcul exact ou approché, puis phrase de réponse. La réponse modèle figure dans l'édition enseignant.
\end{frame}
\begin{frame}{Taux d'évolution moyen}
\textbf{Conduite en 1 h :} 7 min d'automatismes ; 10 min de mise en commun et trace ; 35 min d'exercices essentiels ; 5 min de correction ; 3 min de billet de sortie.

\textbf{Conduite en 2 h :} ajouter 20 min de recherche de l'activité, 20 min d'institutionnalisation, un second exercice autonome et 10 min de correction comparée.

\textbf{Différenciation :} fournir aux élèves fragiles une fiche à quatre cases (données, modèle, calcul, conclusion). Demander aux élèves rapides de produire un contre-exemple, une justification ou une seconde méthode.
\end{frame}
\begin{frame}{Lecture de représentations graphiques}
\begin{itemize}
\item Lire images et antécédents.
\item Comparer deux modèles.
\item Estimer graphiquement.

\end{itemize}
\end{frame}
\begin{frame}{Lecture de représentations graphiques}
\begin{columns}[T,onlytextwidth]
\begin{column}{.57\textwidth}
Deux offres commerciales sont modélisées par \(f(x)=120\times1{,}08^x\) et \(g(x)=180\times0{,}96^x\). Tracer les courbes sur \([0,12]\), lire leur point d'intersection et interpréter les zones où chaque valeur est la plus grande.
\end{column}
\begin{column}{.4\textwidth}
\centering
\includegraphics[width=\linewidth,height=4.5cm,keepaspectratio]{figures/stmg-02-4.pdf}
\end{column}
\end{columns}
\end{frame}
\begin{frame}{Lecture de représentations graphiques}
Lire \(f(x)\) consiste à partir de \(x\) vers la courbe puis vers l'axe des ordonnées. Chercher un antécédent de \(y\) consiste à partir de \(y\) vers la courbe puis vers l'axe des abscisses.
\end{frame}
\begin{frame}{Lecture de représentations graphiques}
Reprendre les données de l'activité et rédiger la solution dans l'ordre : définition des objets, propriété utilisée, calcul exact ou approché, puis phrase de réponse. La réponse modèle figure dans l'édition enseignant.
\end{frame}
\begin{frame}{Lecture de représentations graphiques}
\textbf{Conduite en 1 h :} 7 min d'automatismes ; 10 min de mise en commun et trace ; 35 min d'exercices essentiels ; 5 min de correction ; 3 min de billet de sortie.

\textbf{Conduite en 2 h :} ajouter 20 min de recherche de l'activité, 20 min d'institutionnalisation, un second exercice autonome et 10 min de correction comparée.

\textbf{Différenciation :} fournir aux élèves fragiles une fiche à quatre cases (données, modèle, calcul, conclusion). Demander aux élèves rapides de produire un contre-exemple, une justification ou une seconde méthode.
\end{frame}
\begin{frame}{Équations exponentielles simples}
\begin{itemize}
\item Résoudre avec une même base.
\item Approcher une solution.
\item Vérifier la vraisemblance.

\end{itemize}
\end{frame}
\begin{frame}{Équations exponentielles simples}
Résoudre \(4^x=64\), \(9^x=\dfrac13\) et \(1{,}12^x=2\) : d'abord par même base lorsque c'est possible, puis numériquement pour la dernière. Encadrer la dernière solution entre deux entiers.
\end{frame}
\begin{frame}{Équations exponentielles simples}
Pour \(a>0\) et \(a\neq1\), pour résoudre \(a^x=b\) : utiliser une même base quand c'est possible ; sinon encadrer ou utiliser un outil numérique. Vérifier que \(b>0\).
\end{frame}
\begin{frame}{Équations exponentielles simples}
Reprendre les données de l'activité et rédiger la solution dans l'ordre : définition des objets, propriété utilisée, calcul exact ou approché, puis phrase de réponse. La réponse modèle figure dans l'édition enseignant.
\end{frame}
\begin{frame}{Équations exponentielles simples}
\textbf{Conduite en 1 h :} 7 min d'automatismes ; 10 min de mise en commun et trace ; 35 min d'exercices essentiels ; 5 min de correction ; 3 min de billet de sortie.

\textbf{Conduite en 2 h :} ajouter 20 min de recherche de l'activité, 20 min d'institutionnalisation, un second exercice autonome et 10 min de correction comparée.

\textbf{Différenciation :} fournir aux élèves fragiles une fiche à quatre cases (données, modèle, calcul, conclusion). Demander aux élèves rapides de produire un contre-exemple, une justification ou une seconde méthode.
\end{frame}
\begin{frame}{Modélisation par une fonction exponentielle}
\begin{itemize}
\item Construire un modèle exponentiel.
\item Identifier les paramètres.
\item Prévoir et interpréter.

\end{itemize}
\end{frame}
\begin{frame}{Modélisation par une fonction exponentielle}
\begin{columns}[T,onlytextwidth]
\begin{column}{.57\textwidth}
Une application compte \(12\,000\) utilisateurs au lancement et sa fréquentation augmente de \(9\,\%\) par mois. Construire \(f(t)\), calculer \(f(6)\), puis déterminer numériquement quand elle dépassera \(25\,000\) utilisateurs.
\end{column}
\begin{column}{.4\textwidth}
\centering
\includegraphics[width=\linewidth,height=4.5cm,keepaspectratio]{figures/stmg-02-5.pdf}
\end{column}
\end{columns}
\end{frame}
\begin{frame}{Modélisation par une fonction exponentielle}
Une évolution relative constante de coefficient \(q\) se modélise par \(f(t)=V_0q^t\). \(V_0\) est la valeur initiale et \(q\) le coefficient par unité de temps.
\end{frame}
\begin{frame}{Modélisation par une fonction exponentielle}
Reprendre les données de l'activité et rédiger la solution dans l'ordre : définition des objets, propriété utilisée, calcul exact ou approché, puis phrase de réponse. La réponse modèle figure dans l'édition enseignant.
\end{frame}
\begin{frame}{Modélisation par une fonction exponentielle}
\textbf{Conduite en 1 h :} 7 min d'automatismes ; 10 min de mise en commun et trace ; 35 min d'exercices essentiels ; 5 min de correction ; 3 min de billet de sortie.

\textbf{Conduite en 2 h :} ajouter 20 min de recherche de l'activité, 20 min d'institutionnalisation, un second exercice autonome et 10 min de correction comparée.

\textbf{Différenciation :} fournir aux élèves fragiles une fiche à quatre cases (données, modèle, calcul, conclusion). Demander aux élèves rapides de produire un contre-exemple, une justification ou une seconde méthode.
\end{frame}
\begin{frame}{Bilan et évaluation du chapitre}
\begin{itemize}
\item Mobiliser les acquis du chapitre.
\item Rédiger une modélisation.
\item Corriger une erreur.

\end{itemize}
\end{frame}
\begin{frame}{Bilan et évaluation du chapitre}
À partir de quatre productions d'élèves, repérer et corriger : confusion taux/coefficient, propriété de puissances fausse, erreur de sens de variation et conclusion sans unité.
\end{frame}
\begin{frame}{Bilan et évaluation du chapitre (2)}
\begin{enumerate}
\setcounter{enumi}{0}
\item Une population de \(2\,000\) habitants augmente de \(3\,\%\) par an. Léa écrit : « \(P(n)=2\,000\times3^n\). »

\item Karim écrit : « \(2^x+2^3=2^{x+3}\). »

\item Awa écrit : « \(f(x)=-3\times0{,}5^x\) est décroissante, car \(0{,}5<1\). »

\item Une machine achetée \(12\,000\) euros perd \(8\,\%\) de sa valeur par an. Tojo écrit : « \(V(n)=12\,000\times0{,}92^n\), donc \(V(5)\approx7\,909\). Réponse : \(7\,909\). »

\end{enumerate}
\end{frame}
\begin{frame}{Bilan et évaluation du chapitre}
Pour valider un modèle exponentiel : identifier \(V_0\) et \(q\), calculer avec les unités, vérifier l'ordre de grandeur et interpréter le résultat.
\end{frame}
\begin{frame}{Bilan et évaluation du chapitre}
Reprendre les données de l'activité et rédiger la solution dans l'ordre : définition des objets, propriété utilisée, calcul exact ou approché, puis phrase de réponse. La réponse modèle figure dans l'édition enseignant.
\end{frame}
\begin{frame}{Bilan et évaluation du chapitre}
\textbf{Conduite en 1 h :} 7 min d'automatismes ; 10 min de mise en commun et trace ; 35 min d'exercices essentiels ; 5 min de correction ; 3 min de billet de sortie.

\textbf{Conduite en 2 h :} ajouter 20 min de recherche de l'activité, 20 min d'institutionnalisation, un second exercice autonome et 10 min de correction comparée.

\textbf{Différenciation :} fournir aux élèves fragiles une fiche à quatre cases (données, modèle, calcul, conclusion). Demander aux élèves rapides de produire un contre-exemple, une justification ou une seconde méthode.
\end{frame}
\end{document}
